\section{Benchmark Construction}
\label{sec:benchmark}

The \ours benchmark is designed to evaluate continual graph learning methods on a real, evolving biomedical knowledge graph.
We describe the temporal snapshot construction, temporal difference analysis, continual task sequence design, evaluation tasks, and multimodal feature extraction.

\subsection{Temporal Snapshots}
\label{sec:snapshots}

We construct two temporal snapshots of PrimeKG~\citep{Chandak2023} by querying its upstream biomedical databases at different time points:

\paragraph{Snapshot $t_0$ (June 2021).}
This corresponds to the original PrimeKG release, downloaded from the Harvard Dataverse.
It contains 129,375 nodes across 10 entity types and 8,100,498 edges spanning 30 relation types, integrating data from 20+ biomedical databases including DrugBank, DisGeNET, Gene Ontology, HPO, Reactome, and MONDO.

\paragraph{Snapshot $t_1$ (July 2023).}
We reconstruct PrimeKG by re-querying nine freely accessible upstream databases: Bgee (gene expression), CTD (chemical-disease associations), Gene Ontology, Gene2GO (gene-GO annotations), HPO (phenotype ontology), HPOA (phenotype-disease annotations), MONDO (disease ontology), Uberon (anatomy ontology), and HGNC (gene nomenclature).
Seven databases with restrictive licensing (DrugBank, UMLS, DrugCentral, SIDER) or API changes (DisGeNET) were not re-queried; for these, we carry forward the $t_0$ edges to ensure the $t_1$ graph is a superset of $t_0$ for the corresponding relation types.
The resulting $t_1$ snapshot contains 134,211 nodes and 13,041,729 edges across 25 relation types.

\Cref{tab:dataset_stats} summarizes the statistics of both snapshots.

\begin{table}[t]
\centering
\caption{Dataset statistics for the two temporal snapshots in \ours. The temporal difference reveals substantial real-world knowledge evolution between June 2021 and July 2023.}
\label{tab:dataset_stats}
\small
\begin{tabular}{lcccc}
\toprule
\textbf{Statistic} & \textbf{$t_0$ (June 2021)} & \textbf{$t_1$ (July 2023)} & \textbf{Difference} \\
\midrule
Nodes & 129,375 & 134,211 & +4,836 \\
Edges & 8,100,498 & 13,041,729 & +4,941,231 \\
Node types & 10 & 10 & 0 \\
Relation types & 30 & 25 & $-5$ \\
\midrule
Added edges & -- & -- & +5,730,243 \\
Removed edges & -- & -- & $-889,012$ \\
Persistent edges & -- & -- & 7,211,486 \\
\midrule
Continual tasks & \multicolumn{3}{c}{10 (entity-type grouping)} \\
Evaluation tracks & \multicolumn{3}{c}{LP, KGQA, Node Classification} \\
\bottomrule
\end{tabular}
\end{table}


\subsection{Temporal Difference}
\label{sec:temporal_diff}

The temporal difference between $t_0$ and $t_1$ reveals substantial real-world knowledge evolution:
\begin{itemize}
    \item \textbf{Added edges:} 5,730,243 new triples appear in $t_1$ that were absent in $t_0$, reflecting newly discovered associations (e.g., new gene-GO annotations, updated disease-phenotype links).
    \item \textbf{Removed edges:} 889,012 triples present in $t_0$ are absent from $t_1$, representing deprecated or corrected knowledge (e.g., retracted gene-disease associations, updated ontology hierarchies).
    \item \textbf{Persistent edges:} 7,211,486 triples remain unchanged across both snapshots, providing a stable knowledge backbone.
\end{itemize}

This temporal difference is fundamentally different from synthetic CGL benchmarks that randomly partition a static graph.
The additions and removals follow domain-specific patterns: Gene Ontology annotations account for the largest share of added edges (reflecting the rapid expansion of functional annotations), while disease-phenotype associations show high turnover as HPO and MONDO ontologies are refined.

\subsection{Continual Task Sequence}
\label{sec:task_sequence}

We organize the continual learning task sequence using an \emph{entity-type grouping} strategy.
Rather than randomly partitioning triples into artificial time steps (as in LKGE~\citep{Cui2023}), we group edges by their head and tail entity types to create semantically coherent tasks.
Specifically, we define 10 continual tasks based on the dominant entity-type pairs in PrimeKG:
\begin{enumerate}
    \item Gene/Protein $\leftrightarrow$ Gene/Protein (protein-protein interactions)
    \item Gene/Protein $\leftrightarrow$ Biological Process
    \item Gene/Protein $\leftrightarrow$ Molecular Function
    \item Gene/Protein $\leftrightarrow$ Cellular Component
    \item Gene/Protein $\leftrightarrow$ Disease
    \item Gene/Protein $\leftrightarrow$ Pathway
    \item Disease $\leftrightarrow$ Phenotype
    \item Drug $\leftrightarrow$ Disease (drug indications and contraindications)
    \item Drug $\leftrightarrow$ Side Effect
    \item Anatomy $\leftrightarrow$ Gene/Protein (tissue expression)
\end{enumerate}

Each task is presented sequentially at both $t_0$ and $t_1$, requiring the model to first learn the task on $t_0$ data and then adapt to $t_1$ data while retaining knowledge of previously learned tasks.
This entity-type grouping produces tasks of varying difficulty and size, mirroring the heterogeneous nature of biomedical knowledge evolution.

\subsection{Evaluation Tasks}
\label{sec:eval_tasks}

\ours supports three evaluation tracks:

\paragraph{Link Prediction (Primary).}
Given a query $(h, r, ?)$ or $(?, r, t)$, the model ranks all candidate entities.
We evaluate using Mean Reciprocal Rank (MRR), Hits@$K$ ($K \in \{1, 3, 10\}$), and Area Under the Precision-Recall Curve (AUPRC).
In the continual setting, we additionally report Average Performance (AP), Average Forgetting (AF), Backward Transfer (BWT), and Remembering (REM), computed across all tasks after training on each new task:
\begin{align}
    \text{AP} &= \frac{1}{T} \sum_{i=1}^{T} a_{T,i}, \label{eq:ap} \\
    \text{AF} &= \frac{1}{T-1} \sum_{i=1}^{T-1} \max_{j \in \{1,\ldots,T-1\}} (a_{j,i} - a_{T,i}), \label{eq:af} \\
    \text{BWT} &= \frac{1}{T-1} \sum_{i=1}^{T-1} (a_{T,i} - a_{i,i}), \label{eq:bwt} \\
    \text{REM} &= 1 - \text{AF}, \label{eq:rem}
\end{align}
where $a_{j,i}$ denotes the performance on task $i$ after training on task $j$, and $T$ is the total number of tasks.

\paragraph{Knowledge Graph Question Answering (KGQA).}
We evaluate continual KGQA using a retrieval-augmented generation (RAG) pipeline~\citep{Lewis2020} that retrieves relevant subgraph context from the current KG snapshot and feeds it to an LLM for answer generation.
We measure Exact Match (EM), token-level F1, and accuracy.

\paragraph{Node Classification.}
We classify disease nodes into disease categories using node features derived from the KG structure and multimodal attributes.
We report Macro-F1 and accuracy, along with CL metrics (AP, AF, BWT).

\subsection{Multimodal Features}
\label{sec:features}

A distinctive feature of \ours is the availability of multimodal node representations:

\paragraph{Textual features.}
We encode entity descriptions using BiomedBERT~\citep{Gu2021}, a domain-specific language model pre-trained on PubMed abstracts.
Each node's textual description (e.g., gene function summary, disease definition, drug mechanism of action) is encoded into a 768-dimensional vector using the [CLS] token representation.

\paragraph{Molecular features.}
For drug nodes with available SMILES strings, we compute 2048-bit Morgan fingerprints~\citep{Rogers2004} (radius 2), capturing circular substructure information relevant to molecular similarity and activity prediction.
These binary fingerprints are projected to a dense representation via a learned MLP.

\paragraph{Structural features.}
We extract structural node embeddings using a Relational Graph Convolutional Network (R-GCN)~\citep{Schlichtkrull2018} that captures multi-relational neighborhood information.
The R-GCN produces type-aware embeddings that reflect each node's local graph context.

These multimodal features enable the study of modality-specific forgetting patterns---a novel research direction that is impossible on existing CGL benchmarks lacking multimodal information.
